\documentclass{article}
\usepackage{enumitem}
\usepackage{amsmath}
\begin{document}

\title{Chapter 19: Excretory Products and their Elimination}
\date{}
\maketitle

\begin{enumerate}[label=Q\arabic*.]

\item The glomerular filtration rate (GFR) in a healthy adult is approximately:
\\(A) 125 mL/min (180 L/day)
\\(B) 1.5 mL/min (2 L/day, urine output)
\\(C) 500 mL/min
\\(D) 60 mL/min

\item The process of ultrafiltration in the glomerulus is driven by:
\\(A) Osmotic pressure of plasma proteins
\\(B) Net filtration pressure (glomerular hydrostatic pressure minus capsular and oncotic pressures)
\\(C) Active transport of water
\\(D) Tubular secretion pressure

\item Which of the following substances is freely filtered at the glomerulus and completely reabsorbed in the proximal tubule, making it useful for measuring GFR?
\\(A) Creatinine
\\(B) Inulin
\\(C) Urea
\\(D) Para-aminohippuric acid (PAH)

\item ADH (vasopressin) increases water reabsorption by:
\\(A) Opening Na$^+$ channels in the proximal tubule
\\(B) Inserting aquaporin-2 (AQP2) channels into the apical membrane of collecting duct cells
\\(C) Stimulating Na$^+$/K$^+$ ATPase in distal tubule
\\(D) Constricting afferent arterioles

\item Assertion: Aldosterone acts on the distal convoluted tubule and collecting duct.\\
Reason: Aldosterone stimulates synthesis of Na$^+$ channels and Na$^+$/K$^+$ ATPase, increasing Na$^+$ reabsorption and K$^+$ secretion.
\\(A) Both true, Reason is correct explanation
\\(B) Both true, Reason is not correct explanation
\\(C) Assertion true, Reason false
\\(D) Assertion false, Reason true

\item The countercurrent multiplier in the loop of Henle creates a:
\\(A) Concentration gradient in the medullary interstitium that drives water reabsorption
\\(B) Pressure gradient that drives filtration
\\(C) Blood flow gradient from cortex to medulla
\\(D) Osmotic gradient that drives Na$^+$ secretion

\item Uricotelism (excretion of uric acid) is found in:
\\(A) Aquatic animals (fish, tadpoles)
\\(B) Terrestrial animals needing to conserve water (reptiles, birds, insects)
\\(C) Mammals including humans
\\(D) Marine invertebrates

\item The juxtaglomerular apparatus (JGA) is important because it:
\\(A) Produces erythropoietin
\\(B) Regulates GFR by secreting renin in response to low blood pressure
\\(C) Reabsorbs glucose in proximal tubule
\\(D) Secretes ADH directly

\item Renal threshold for glucose is approximately:
\\(A) 70 mg/dL
\\(B) 180 mg/dL
\\(C) 300 mg/dL
\\(D) 90 mg/dL

\item In the RAAS (Renin-Angiotensin-Aldosterone System), the correct sequence is:
\\(A) Renin $\rightarrow$ Angiotensinogen $\rightarrow$ Angiotensin I $\rightarrow$ Angiotensin II $\rightarrow$ Aldosterone
\\(B) Aldosterone $\rightarrow$ Renin $\rightarrow$ Angiotensin I $\rightarrow$ Angiotensin II
\\(C) Renin $\rightarrow$ Angiotensin II $\rightarrow$ Angiotensin I $\rightarrow$ Aldosterone
\\(D) Angiotensinogen $\rightarrow$ Aldosterone $\rightarrow$ Renin $\rightarrow$ Angiotensin II

\item The proximal convoluted tubule (PCT) is responsible for reabsorption of approximately what percentage of filtered load?
\\(A) 25\%
\\(B) 50\%
\\(C) 65--70\% (of Na$^+$, water, and other solutes)
\\(D) 90\%

\item Hemodialysis mimics which kidney function?
\\(A) Hormone production (EPO)
\\(B) Filtration of blood and removal of waste products/excess fluid across a semipermeable membrane
\\(C) Active secretion of drugs
\\(D) Regulation of blood pressure via RAAS

\item ANP (Atrial Natriuretic Peptide) is secreted when:
\\(A) Blood pressure falls
\\(B) Atria are stretched due to increased blood volume/pressure; ANP promotes Na$^+$ and water excretion
\\(C) Aldosterone levels are high
\\(D) Renin is released from JGA

\item Match the following nephron segments with their primary functions:\\
1. PCT $\rightarrow$ (a) Obligatory reabsorption of 65--70\% Na$^+$, glucose, amino acids\\
2. Loop of Henle $\rightarrow$ (b) Countercurrent multiplication to create medullary gradient\\
3. DCT $\rightarrow$ (c) Aldosterone-regulated Na$^+$, K$^+$ balance\\
4. Collecting duct $\rightarrow$ (d) ADH-regulated water reabsorption
\\(A) 1-a, 2-b, 3-c, 4-d
\\(B) 1-b, 2-a, 3-d, 4-c
\\(C) 1-c, 2-d, 3-a, 4-b
\\(D) 1-d, 2-c, 3-b, 4-a

\item Uremia is a condition characterized by:
\\(A) Excess urine production
\\(B) Accumulation of urea and other nitrogenous wastes in blood due to renal failure
\\(C) Urine glucose above normal
\\(D) Deficiency of ADH

\end{enumerate}
\end{document}
